\begin{longtable}[c]{|m{0.25\textwidth}|m{0.7\textwidth}|}
    \caption{Tareas Específicas para medición de la Usabilidad del \textit{middleware}.}
    \label{tab:tareasEspecificas} \\
    
     \hline
    \rowcolor{gray!50}
    \multicolumn{1}{|c|}{\textbf{Tarea}} & \multicolumn{1}{c|}{\textbf{Pasos}} \\
    \hline
    \endfirsthead

    \hline
    \rowcolor{gray!50}
    \multicolumn{1}{|c|}{\textbf{Tarea}} & \multicolumn{1}{c|}{\textbf{Pasos}} \\
    \hline
    \endhead

    \multicolumn{2}{r}{\textit{Continúa en la siguiente página...}} \\
    \endfoot
    
    \endlastfoot

    Tarea 01 & \begin{itemize}
        \item Dirigirse a la página de inicio de sesión.
        \item El usuario debe tener una conexión a Neo4j almacenada y activa.
    \end{itemize} \\
    \hline
    Tarea 02 & \begin{itemize}
        \item Acceder a la pantalla de traducción de consultas.
        \item Seleccionar una conexión a Neo4j activa.
        \item Ingresar una sentencia SQL \texttt{INSERT} válida.
        \item Dar clic en el boton \textbf{Traducir}.
    \end{itemize} \\
    \hline
    Tarea 03 & \begin{itemize}
        \item Acceder a la pantalla de traducción de consultas.
        \item Seleccionar una conexión a Neo4j activa.
        \item Ingresar una sentencia SQL \texttt{INSERT} válida.
        \item Dar clic en el boton \textbf{Traducir}.
    \end{itemize} \\
    \hline
    Tarea 04 & \begin{itemize}
        \item Acceder a la pantalla de traducción de consultas.
        \item Seleccionar una conexión a Neo4j activa.
        \item Ingresar una sentencia SQL \texttt{INSERT} válida.
        \item Dar clic en el boton \textbf{Traducir}.
    \end{itemize} \\
    \hline
    Tarea 05 & \begin{itemize}
        \item Acceder a la pantalla de traducción de consultas.
        \item Seleccionar una conexión a Neo4j activa.
        \item Ingresar una sentencia SQL \texttt{INSERT} válida.
        \item Dar clic en el boton \textbf{Traducir}.
    \end{itemize} \\
    \hline
\end{longtable}